% ============================================================
% Conclusions (target: Computers & Security)
% ============================================================
\section{Conclusions}
\label{sec:conclusion}

Modern phishing detection works in aggregate, yet a residual fraction, spear phishing and BEC, remains the most
costly and almost invisible to content filters. We have argued that this is a structural property: a targeted
attack succeeds by respecting the victim's domain knowledge, so the maliciousness signal lies not in the content
but in its inconsistency with the graph of that knowledge and in the dynamics of its violation.

\paragraph{Answers to the research questions.}
\begin{description}[noitemsep,topsep=2pt]
  \item[RQ1 (static signal, T1):] Yes, conditionally. The cross-layer inconsistency invariant beats a single
  contact graph and removes false alarms from legitimate new senders (Section~\ref{sec:res-inv1}), provided there
  is a shared, observable domain-knowledge edge, and the gain vanishes under the layer/rhythm shuffle control.
  \item[RQ2 (temporal signal, T2):] Yes. The temporal invariant removes the blindness of the static layers and,
  unlike a volume-based detector, does not collapse against a stealthy attacker (Section~\ref{sec:res-inv2}); the
  advantage is causally temporal and transfers to three real graphs.
  \item[RQ3 (measurement, T3):] Measurable, and governed by relation rather than rank. Cold external impersonation
  leaves a margin of $\mathrm{Susc}\!\approx\!0.96$ at every level, whereas in the lateral channel it is the
  periphery that is most exposed, not the executives (Section~\ref{sec:susceptibility}). At a realistic base rate
  the detectors act as a triage instrument, with effectiveness defined operationally through low-FPR sensitivity
  under leak controls and a predictivity probe (Section~\ref{sec:res-rigor}).
\end{description}

\paragraph{Contributions.}
\begin{enumerate}[noitemsep,topsep=2pt]
  \item The concept and formalization of the domain knowledge graph, with the thesis that spear phishing is its
  violation, uniformly for lateral/BEC and external impersonation.
  \item Two complementary detection invariants (cross-layer inconsistency and cascade dynamics), with causal
  controls and validation on three real communication graphs.
  \item A measured organizational susceptibility metric showing that the residual margin is governed by relation,
  and that in the lateral channel the periphery is more exposed than the top.
  \item A leak-aware evaluation methodology and an operational definition of spear-phishing detection
  effectiveness.
\end{enumerate}

\paragraph{Limitations (summary).} The maliciousness is synthetic, so we validate the mechanism rather than the
detection of real intrusions; the susceptibility measurement inherits the assumptions of the event model; and
scale and the operational threshold remain limitations. Section~\ref{sec:discussion} gives the details.

The contribution is a concept, a method, and a measurement discipline, rather than a finished product. Spear
phishing is a violation of domain knowledge, and the graph of that knowledge makes both the attack and the
susceptibility to it measurable. This points toward a defense that asks not whether a message looks malicious but
whether it is consistent with what the organization knows about itself.
